\section{SA-HEACS architecture}\label{sec:arch}

SA-HEACS has three components joined by one attested request path: the client shell, the
orchestrator and the secure clusters (\cref{fig:arch}).

\begin{figure}[htbp]
\centering
% Figure 1: SA-HEACS architecture. Everything in the dashed client box is attacker-readable.
\begin{tikzpicture}[
    font=\sffamily\footnotesize,
    box/.style={draw=boxgray, rounded corners=2pt, fill=white, align=left,
                inner sep=5pt, text width=#1},
    box/.default=3.9cm,
    core/.style={box=3.9cm, draw=accent, line width=0.9pt, fill=accentfill, minimum height=11mm},
    region/.style={rounded corners=4pt, inner sep=7pt},
    title/.style={font=\sffamily\footnotesize\bfseries, anchor=south west, inner sep=0pt},
    arr/.style={-{Stealth[length=5pt]}, draw=boxgray, line width=0.6pt},
    lbl/.style={font=\sffamily\scriptsize, text=textgray, align=left}]

  % cloud side (laid out first; the client row is placed above it)
  \node[box] (edge) {\textbf{Edge gateway}\\[1pt]
      \textcolor{textgray}{Verify attestation, nonce}\\\textcolor{textgray}{Rate limits, query monitor}};
  \node[box, right=5mm of edge] (orch) {\textbf{AI orchestrator}\\[1pt]
      \textcolor{textgray}{Tier check, locality routing}\\\textcolor{textgray}{Micro-batching, prefetch}};
  \node[box, right=5mm of orch] (life) {\textbf{Lifecycle pipeline}\\[1pt]
      \textcolor{textgray}{Scores IP value and cost}\\\textcolor{textgray}{Sets each feature's tier}};
  \node[core, below=9mm of edge] (gpu) {\textbf{Tensor / GPU nodes}\\[1pt]\textcolor{textgray}{Neural filters, model weights}};
  \node[core, below=9mm of orch] (vec) {\textbf{Vector math nodes}\\[1pt]\textcolor{textgray}{Geometry, transforms}};
  \node[core, below=9mm of life] (ren) {\textbf{Rendering nodes}\\[1pt]\textcolor{textgray}{Colour science, renders}};
  \node[lbl, below=2.5mm of vec, align=center] (seg) {Segmented clusters: \tier{cloud\_only} code and weights never leave this boundary};
  \node[title] (stitle) at ($(edge.north west)+(0,2.5mm)$) {Vendor cloud (trusted; protected IP lives only here)};
  \begin{scope}[on background layer]
    \node[region, draw=boxgray, fit=(edge)(life)(gpu)(ren)(seg)(stitle)] (cloud) {};
  \end{scope}

  % client side
  \node[box=6.6cm, above=27mm of edge.north west, anchor=south west] (shell) {\textbf{Client shell}\\[1pt]
      \textcolor{textgray}{UI, canvas navigation, brush strokes, files}\\
      \textcolor{textgray}{\tier{client\_ok} features; attested device key}};
  \node[box=5.2cm, dashed, anchor=north east] (recover) at (life.north east |- shell.north)
      {\textbf{What an attacker can recover}\\[1pt]
      \textcolor{textgray}{Shell code, API schema, own inputs and outputs}\\
      \textcolor{textgray}{Not cloud-only code or weights}};
  \node[title] (ctitle) at ($(shell.north west)+(0,2.5mm)$) {Client device (untrusted, attacker-controlled)};
  \begin{scope}[on background layer]
    \node[region, draw=boxgray, dashed, fit=(shell)(recover)(ctitle)] (client) {};
  \end{scope}

  % request path: vertical arrows between the shell and the edge gateway
  \coordinate (dn) at ($(edge.north west)+(2.9,0)$);
  \coordinate (up) at ($(edge.north west)+(1.0,0)$);
  \draw[arr] (dn |- shell.south) -- node[lbl, right=3pt, pos=0.42]
      {Down: attested, encrypted request (\cref{alg:rpc})\\Up: outputs only, watermarked; prefetched results} (dn);
  \draw[arr] (up) -- (up |- shell.south);
  \draw[arr] (edge) -- (orch);
  \draw[arr] (life) -- (orch);
  \draw[arr] (orch.south) -- ++(0,-4.5mm) -| (gpu.north);
  \draw[arr] (orch.south) -- (vec.north);
  \draw[arr] (orch.south) -- ++(0,-4.5mm) -| (ren.north);
\end{tikzpicture}

\caption{SA-HEACS architecture: client shell, attested RPC, orchestrator and segmented
clusters. Everything inside the dashed client boundary is assumed readable by the attacker;
the blue clusters hold the protected IP and return only outputs.}
\label{fig:arch}
\end{figure}

\subsection{Client shell (the untrusted boundary)}
The shell handles input, canvas navigation, local brush strokes, file handling, UI
rendering and any feature in the \tier{client\_ok} tier. It holds no cloud-only code or
weights. Its only secrets are a device identity bound to platform attestation and
short-lived session keys. We assume it will be fully decompiled and cloned; the design goal
is that a perfect clone gains a UI and commodity features, not the protected capabilities.

\subsection{Secure clusters (the protected core)}
Cloud-only features run as segmented microservices grouped by hardware need: tensor and GPU
inference, vector and geometry maths, and rendering or colour-science pipelines. Each
feature runs under its own service identity with least-privilege access to its weights.
Clients never receive code, weights or intermediate tensors that would let a substitute be
rebuilt more cheaply than by black-box queries; responses carry only the final output at the
resolution the product needs.

\subsection{Orchestrator (the control plane)}
The orchestrator sits between the edge and the clusters. It enforces each feature's tier,
routes \tier{hybrid} and \tier{client\_ok} work to the cheaper side, groups compatible
requests into micro-batches, and prefetches likely next operations. \Cref{sec:orch} details
it.

\subsection{Request protocol}
Every remote call carries an attestation-bound signature, a fresh nonce and an encrypted
body (\cref{alg:rpc}). The client derives a per-request key from an ephemeral X25519
exchange \cite{bernstein2006curve25519,rfc7748} with the server's static key, expanded with
HKDF \cite{rfc5869} using the nonce as salt, and encrypts with AES-GCM \cite{nist80038d}.
The ephemeral public key, nonce and timestamp are signed with an Ed25519 device key
\cite{bernstein2012ed25519,rfc8032} that, in production, is bound to a platform attestation
verdict, and the same bytes are passed to AES-GCM as associated data. The server rejects
stale or replayed requests and unattested devices before any protected code runs. Request
and response use disjoint 12-byte GCM nonces under the per-request key.

\begin{algorithm}[htbp]
\caption{Attested invocation of a cloud-only feature}\label{alg:rpc}
\small
\begin{algorithmic}[1]
\Statex \textbf{Client:}
\State \((e, E) \gets \Kw{X25519.Generate}()\)
\State \(n \gets \Kw{Random}(16)\); \(t \gets \Kw{Now}()\)
\State \(\sigma \gets \Kw{Sign}(K_{\text{device}},\, E \parallel n \parallel t)\) \Comment{\(K_{\text{device}}\) bound to platform attestation}
\State \(k \gets \Kw{HKDF}(\Kw{ECDH}(e, S_{\text{pub}}),\ \text{salt}=n)\)
\State \(c \gets \Kw{AES-GCM.Enc}(k,\ n[0{:}12],\ \{\text{feature}, \text{args}\},\ \text{aad}=E \parallel n \parallel t)\)
\State \textbf{send} \((\text{device\_id}, E, n, t, \sigma, c)\)
\Statex \textbf{Server:}
\If{\(|\Kw{Now}() - t| > \Delta\) \textbf{or} \(n \in \text{SeenNonces}\)} \textbf{reject} \EndIf
\If{\textbf{not} \(\Kw{Verify}(\text{AttestedKey}[\text{device\_id}], \sigma)\)} \textbf{reject} \EndIf
\If{\textbf{not} \(\Kw{RateLimiter.Allow}(\text{account}, \text{feature})\)} \textbf{reject} \Comment{\cref{sec:security}} \EndIf
\State \(k \gets \Kw{HKDF}(\Kw{ECDH}(S_{\text{priv}}, E),\ \text{salt}=n)\)
\State \(\text{req} \gets \Kw{AES-GCM.Dec}(k,\ n[0{:}12],\ c,\ \text{aad})\)
\State \(y \gets \Kw{Orchestrator.Submit}(\text{req.feature}, \text{req.args})\) \Comment{\cref{sec:orch}}
\State \Return \(\Kw{AES-GCM.Enc}(k,\ n[12{:}16] \parallel 0^{8},\ \Kw{Watermark}(y),\ \text{aad})\)
\end{algorithmic}
\end{algorithm}

The runnable version, \code{attested\_rpc.py} (\cref{app:rpc}), implements lines 1--8 and
10--13, without the rate limiter and watermark, using the Python \code{cryptography} package
\cite{pyca_cryptography} and FastAPI \cite{fastapi}. Running it shows a valid call returning
the filtered result and a cloned client with an unattested key rejected with HTTP~401. In
production, line~8 would consume a Play Integrity or App Attest verdict
\cite{google_playintegrity,apple_appattest} instead of a locally registered Ed25519 key, and
the nonce cache would be bounded by the skew window \(\Delta\).
